% !TEX program = xelatex
% Options for packages loaded elsewhere
\PassOptionsToPackage{unicode}{hyperref}
\PassOptionsToPackage{hyphens}{url}
\PassOptionsToPackage{dvipsnames,svgnames,x11names}{xcolor}
\documentclass[
  8pt,
  twocolumn]{extarticle}
\usepackage{xcolor}
\usepackage[landscape,margin=0.38in]{geometry}
\usepackage{amsmath,amssymb}
\setcounter{secnumdepth}{-\maxdimen} % remove section numbering
\usepackage{iftex}
\ifPDFTeX
  \usepackage[T1]{fontenc}
  \usepackage[utf8]{inputenc}
  \usepackage{textcomp} % provide euro and other symbols
\else % if luatex or xetex
  \usepackage{unicode-math} % this also loads fontspec
  \defaultfontfeatures{Scale=MatchLowercase}
  \defaultfontfeatures[\rmfamily]{Ligatures=TeX,Scale=1}
\fi
\usepackage{lmodern}
\ifPDFTeX\else
  % xetex/luatex font selection
\fi
% Use upquote if available, for straight quotes in verbatim environments
\IfFileExists{upquote.sty}{\usepackage{upquote}}{}
\IfFileExists{microtype.sty}{% use microtype if available
  \usepackage[]{microtype}
  \UseMicrotypeSet[protrusion]{basicmath} % disable protrusion for tt fonts
}{}
\makeatletter
\@ifundefined{KOMAClassName}{% if non-KOMA class
  \IfFileExists{parskip.sty}{%
    \usepackage{parskip}
  }{% else
    \setlength{\parindent}{0pt}
    \setlength{\parskip}{6pt plus 2pt minus 1pt}}
}{% if KOMA class
  \KOMAoptions{parskip=half}}
\makeatother
\setlength{\emergencystretch}{3em} % prevent overfull lines
\providecommand{\tightlist}{%
  \setlength{\itemsep}{0pt}\setlength{\parskip}{0pt}}
\usepackage{microtype}
\usepackage{mathtools}
\usepackage{amssymb}
\setlength{\columnsep}{0.32in}
\setlength{\parindent}{0pt}
\setlength{\parskip}{2pt}
\setlength{\abovedisplayskip}{4pt}
\setlength{\belowdisplayskip}{4pt}
\setlength{\abovedisplayshortskip}{3pt}
\setlength{\belowdisplayshortskip}{3pt}
\allowdisplaybreaks
\sloppy
\emergencystretch=3em
\usepackage{bookmark}
\IfFileExists{xurl.sty}{\usepackage{xurl}}{} % add URL line breaks if available
\urlstyle{same}
\hypersetup{
  colorlinks=true,
  linkcolor={black},
  filecolor={Maroon},
  citecolor={Blue},
  urlcolor={black},
  pdfcreator={LaTeX via pandoc}}

\author{}
\date{}

\begin{document}

\section{Formula Sheet Part 3: Stationarity, AR Models, and
ARCH}\label{formula-sheet-part-3-stationarity-ar-models-and-arch}

This is a two-column topic split from
\texttt{All\_Formulas\_From\_Chatnotes.md}.

\subsection{Conventions}\label{conventions}

\[
t,T,h,k,n,p,q,N
\]

Variables: \(t\) is a time index; \(T\) is sample size or the final
observed time; \(h\) is a forecast horizon; \(k\) is a lag or number of
model parameters depending on context; \(n\) is a return horizon;
\(p,q\) are lag orders; \(N\) is the number of assets.

\[
F_t
\]

Variables: \(F_t\) is the information set available at time \(t\).

Returns may be written as decimals or percentages. Check the scale
before applying VaR or standard deviation formulas.

\subsection{Stationarity, Random Walks, ADF, and AR
Models}\label{stationarity-random-walks-adf-and-ar-models}

\subsubsection{Covariance stationarity}\label{covariance-stationarity}

\[
\begin{aligned}
\mathbb E(Y_t)&=\mu\\
\operatorname{Var}(Y_t)&=\sigma^2\\
\operatorname{Cov}(Y_t,Y_{t-j})&=\gamma_j
\end{aligned}
\]

Variables: stationarity means constant mean, constant variance, and
autocovariance depending only on lag \(j\).

\subsubsection{Stationarity rule of
thumb}\label{stationarity-rule-of-thumb}

\[
\begin{aligned}
\log\text{ prices}&\quad \text{often non-stationary}\\
\log\text{ returns}&\quad \text{often stationary}
\end{aligned}
\]

Variables: this is a modelling guide, not a formal test.

\subsubsection{Random walk}\label{random-walk}

\[
p_t=p_{t-1}+\epsilon_t
\]

Variables: \(p_t\) is usually log price; \(\epsilon_t\) is the shock.

\subsubsection{Random walk successive
substitution}\label{random-walk-successive-substitution}

\[
p_t=p_0+\epsilon_1+\cdots+\epsilon_t
\]

Variables: \(p_0\) is the initial value.

\subsubsection{Random walk moments}\label{random-walk-moments}

\[
\mathbb E(p_t)=p_0,\qquad \operatorname{Var}(p_t)=\sigma^2t
\]

Variables: \(\sigma^2\) is shock variance. If \(p_0=0\), the note's
simplified mean is \(\mathbb E(p_t)=0\).

\subsubsection{Random walk with drift}\label{random-walk-with-drift}

\[
p_t=\mu+p_{t-1}+\epsilon_t
\]

Variables: \(\mu\) is drift.

\subsubsection{Random walk with drift successive
substitution}\label{random-walk-with-drift-successive-substitution}

\[
p_t=p_0+\mu t+\epsilon_1+\cdots+\epsilon_t
\]

Variables: drift accumulates linearly over time.

\subsubsection{Random walk with drift
moments}\label{random-walk-with-drift-moments}

\[
\mathbb E(p_t)=p_0+\mu t,\qquad \operatorname{Var}(p_t)=\sigma^2t
\]

Variables: if \(p_0=0\), the note's simplified mean is \(\mu t\).

\subsubsection{Unit-root AR form}\label{unit-root-ar-form}

\[
y_t=\phi_1y_{t-1}+\text{error}_t
\]

Variables: \(\phi_1\) is the autoregressive coefficient.

\subsubsection{Unit root condition}\label{unit-root-condition}

\[
\phi_1=1
\]

Variables: a unit root implies random-walk-type non-stationarity.

\subsubsection{ADF regression}\label{adf-regression}

\[
\Delta y_t=c_t+\phi_cy_{t-1}+\beta_1\Delta y_{t-1}+\cdots+\beta_p\Delta y_{t-p}+\epsilon_t
\]

Variables: \(\Delta y_t=y_t-y_{t-1}\); \(c_t\) is deterministic
component; \(\phi_c\) is the unit-root test coefficient; \(\beta_j\)
controls lagged differences.

\subsubsection{ADF coefficient
relationship}\label{adf-coefficient-relationship}

\[
\phi_c=\phi_1-1
\]

Variables: \(\phi_c=0\) corresponds to \(\phi_1=1\).

\subsubsection{ADF hypotheses}\label{adf-hypotheses}

\[
\begin{aligned}
H_0&:\phi_c=0\quad \text{unit root / non-stationary}\\
H_1&:\phi_c<0\quad \text{stationary}
\end{aligned}
\]

Variables: the ADF alternative is one-sided and negative.

\subsubsection{ADF statistic}\label{adf-statistic}

\[
\operatorname{ADF}=\frac{\hat\phi_c}{\operatorname{se}(\hat\phi_c)}
\]

Variables: \(\hat\phi_c\) is the estimated ADF coefficient.

\subsubsection{ADF decision rule}\label{adf-decision-rule}

\[
\operatorname{ADF}<\text{ADF critical value}\Rightarrow \text{reject }H_0
\]

Variables: ``less than'' means more negative than the critical value.

\subsubsection{Drift-only ADF deterministic
component}\label{drift-only-adf-deterministic-component}

\[
c_t=c
\]

Variables: \(c\) is a constant drift/intercept.

\subsubsection{Trend ADF deterministic
component}\label{trend-adf-deterministic-component}

\[
c_t=c_0+c_1t
\]

Variables: \(c_0\) is an intercept; \(c_1\) is deterministic trend
slope.

\subsubsection{Martingale difference
sequence}\label{martingale-difference-sequence}

\[
\mathbb E(\epsilon_t\mid F_{t-1})=0
\]

Variables: \(\epsilon_t\) has zero conditional mean given past
information.

\subsubsection{AR(1) model}\label{ar1-model}

\[
y_t=c+\phi_1y_{t-1}+\epsilon_t
\]

Variables: \(c\) is intercept; \(\phi_1\) is the AR(1) coefficient.

\subsubsection{AR(1) conditional mean}\label{ar1-conditional-mean}

\[
\mathbb E(y_t\mid F_{t-1})=c+\phi_1y_{t-1}
\]

Variables: the conditional mean is the one-step forecast given past
information.

\subsubsection{AR(1) unconditional mean}\label{ar1-unconditional-mean}

\[
\mathbb E(y_t)=\frac{c}{1-\phi_1}
\]

Variables: this requires stationarity, typically \(|\phi_1|<1\).

\subsubsection{AR(1) conditional
variance}\label{ar1-conditional-variance}

\[
\operatorname{Var}(\epsilon_t\mid F_{t-1})=\sigma_\epsilon^2
\]

Variables: \(\sigma_\epsilon^2\) is innovation variance.

\[
\operatorname{Var}(y_t\mid F_{t-1})=\sigma_\epsilon^2
\]

Variables: uncertainty in \(y_t\) conditional on \(F_{t-1}\) comes from
the new shock.

\subsubsection{AR(1) unconditional
variance}\label{ar1-unconditional-variance}

\[
\operatorname{Var}(y_t)=\frac{\sigma_\epsilon^2}{1-\phi_1^2}
\]

Variables: this requires stationarity.

\subsubsection{AR(1) one-step forecast}\label{ar1-one-step-forecast}

\[
\mathbb E(y_{t+1}\mid F_t)=c+\phi_1y_t
\]

Variables: \(y_t\) is known at time \(t\).

\subsubsection{AR(1) two-step forecast}\label{ar1-two-step-forecast}

\[
\mathbb E(y_{t+2}\mid F_t)=c+\phi_1\mathbb E(y_{t+1}\mid F_t)=c(1+\phi_1)+\phi_1^2y_t
\]

Variables: the unknown future \(y_{t+1}\) is replaced by its forecast.

\subsubsection{AR(1) long-horizon
forecast}\label{ar1-long-horizon-forecast}

\[
\mathbb E(y_{t+h}\mid F_t)\to \frac{c}{1-\phi_1}
\]

Variables: forecasts mean-revert to the unconditional mean when the
AR(1) is stationary.

\subsection{ARCH and Volatility
Models}\label{arch-and-volatility-models}

\subsubsection{White noise}\label{white-noise}

\[
\begin{aligned}
\mathbb E(u_t)&=0\\
\operatorname{Var}(u_t)&=\sigma^2\\
\operatorname{Corr}(u_t,u_{t-k})&=0,\quad k\ge1
\end{aligned}
\]

Variables: \(u_t\) is a white-noise shock.

\subsubsection{MDS conditional mean}\label{mds-conditional-mean}

\[
\mathbb E(u_t\mid F_{t-1})=0
\]

Variables: an MDS is unpredictable from past information.

\subsubsection{MDS implies zero unconditional
mean}\label{mds-implies-zero-unconditional-mean}

\[
\mathbb E(u_t)=\mathbb E[\mathbb E(u_t\mid F_{t-1})]=0
\]

Variables: this uses the law of iterated expectations.

\subsubsection{MDS implies no
autocovariance}\label{mds-implies-no-autocovariance}

\[
\operatorname{Cov}(u_t,u_{t-k})=0,\quad k\ge1
\]

Variables: \(u_{t-k}\) is known at time \(t-1\).

\subsubsection{AR(p) model}\label{arp-model}

\[
y_t=c+\phi_1y_{t-1}+\cdots+\phi_py_{t-p}+\epsilon_t
\]

Variables: \(p\) is the autoregressive lag order.

\subsubsection{AR(p) unconditional mean}\label{arp-unconditional-mean}

\[
\mathbb E(y_t)=\frac{c}{1-\phi_1-\cdots-\phi_p}
\]

Variables: denominator uses the sum of AR coefficients.

\subsubsection{AR(p) one-step forecast}\label{arp-one-step-forecast}

\[
\mathbb E(y_{t+1}\mid F_t)=c+\phi_1y_t+\phi_2y_{t-1}+\cdots+\phi_py_{t-p+1}
\]

Variables: all lagged \(y\) values on the right are known at time \(t\).

\subsubsection{AR(p) selection}\label{arp-selection}

\[
\text{lowest }\operatorname{AIC}=\text{preferred model}
\]

Variables: \(\operatorname{AIC}\) is Akaike information criterion.

\subsubsection{Mean equation with
residual}\label{mean-equation-with-residual}

\[
r_t=\mathbb E(r_t\mid F_{t-1})+\epsilon_t
\]

Variables: \(\epsilon_t\) is the return shock after using information at
\(t-1\).

\subsubsection{Constant-variance
benchmark}\label{constant-variance-benchmark}

\[
\operatorname{Var}(\epsilon_t\mid F_{t-1})=\sigma_\epsilon^2
\]

Variables: this benchmark has no time-varying conditional volatility.

\subsubsection{ARCH LM mean model}\label{arch-lm-mean-model}

\[
y_t=x_t'\beta+\epsilon_t
\]

Variables: \(x_t\) is a regressor vector; \(\beta\) is a coefficient
vector.

\subsubsection{ARCH LM auxiliary
regression}\label{arch-lm-auxiliary-regression}

\[
\hat\epsilon_t^2=\rho_0+\rho_1\hat\epsilon_{t-1}^2+\cdots+\rho_q\hat\epsilon_{t-q}^2+v_t
\]

Variables: \(\hat\epsilon_t^2\) is squared residual; \(q\) is the number
of ARCH lags tested.

\subsubsection{ARCH LM hypotheses and
statistic}\label{arch-lm-hypotheses-and-statistic}

\[
\begin{aligned}
H_0&:\rho_1=\cdots=\rho_q=0\quad \text{no ARCH effects}\\
H_1&:\text{at least one }\rho_j\ne0\quad \text{ARCH effects exist}\\
TR^2&\sim\chi_q^2
\end{aligned}
\]

Variables: \(R^2\) comes from the auxiliary regression.

\subsubsection{ARCH LM decision rule}\label{arch-lm-decision-rule}

\[
TR^2>\chi_q^2\text{ critical value}\Rightarrow \text{reject }H_0
\]

Variables: rejecting means volatility is time-varying and predictable.

\subsubsection{ARCH model structure}\label{arch-model-structure}

\[
\epsilon_t=\sigma_tu_t,\qquad u_t\sim\mathrm{iid}(0,1)
\]

Variables: \(\sigma_t\) is conditional standard deviation; \(u_t\) is a
standardised shock.

\subsubsection{ARCH conditional moments}\label{arch-conditional-moments}

\[
\mathbb E(\epsilon_t\mid F_{t-1})=0,\qquad \operatorname{Var}(\epsilon_t\mid F_{t-1})=\sigma_t^2
\]

Variables: \(\sigma_t^2\) is conditional variance.

\subsubsection{ARCH(1)}\label{arch1}

\[
\sigma_t^2=\alpha_0+\alpha_1\epsilon_{t-1}^2
\]

Variables: \(\alpha_0\) is variance intercept; \(\alpha_1\) is the ARCH
coefficient.

\subsubsection{ARCH(1) conditions}\label{arch1-conditions}

\[
\alpha_0>0,\qquad \alpha_1\ge0,\qquad \alpha_1<1
\]

Variables: these keep variance positive and finite.

\subsubsection{ARCH(q)}\label{archq}

\[
\sigma_t^2=\alpha_0+\alpha_1\epsilon_{t-1}^2+\cdots+\alpha_q\epsilon_{t-q}^2
\]

Variables: \(q\) is the number of lagged squared shocks.

\subsubsection{ARCH(q) conditions}\label{archq-conditions}

\[
\alpha_0>0,\qquad \alpha_i\ge0,\qquad \sum_{i=1}^{q}\alpha_i<1
\]

Variables: \(\sum\alpha_i\) measures volatility persistence in ARCH(q).

\subsubsection{ARCH(1) unconditional
moments}\label{arch1-unconditional-moments}

\[
\begin{aligned}
\mathbb E(\epsilon_t)&=0\\
\operatorname{Cov}(\epsilon_t,\epsilon_{t-k})&=0,\quad k\ge1\\
\operatorname{Var}(\epsilon_t)&=\frac{\alpha_0}{1-\alpha_1}
\end{aligned}
\]

Variables: ARCH errors are uncorrelated, but squared errors can be
autocorrelated.

\subsubsection{Information criteria}\label{information-criteria}

\[
AIC=-2\ell+2k,\qquad BIC=-2\ell+k\log(T)
\]

Variables: \(\ell\) is log-likelihood; \(k\) is number of estimated
parameters; lower is better.

\end{document}
